\section{Benchmark design and experimental protocol}
\label{sec:benchmark}

The claims this model makes are recovery claims. It claims to find the lagged
parent sets \(\gP_{k,\ell}\) that generated a multiscale log, to track the
latent regime \(S_t^{(k)}\) and the latent factor \(H_t^{(k)}\) that drive it,
and to do both while the counts and durations reconcile exactly across scales.
None of these can be checked against a quantity the data supply, so the
evaluation has to supply the quantity itself. This section states what the
generated benchmark contains, how a scenario is configured, and, before any of
that is used, how a scenario is screened for whether it contains anything to
recover at all.

\subsection{Why recovery is not measurable on an observed log}
\label{sec:benchmark_why}

On an observed event log the true parent sets, the true regime path and the
true parameters \(\Theta^\star\) are all unavailable, and no amount of care in
the fitting procedure produces them. A precision against \(\gP_{k,\ell}^\star\)
is not a hard quantity to compute on such a log; it is an undefined one. What
an observed log does support is prediction on held-out time, exact
reconciliation, the stability of a selected edge across subsamples, and
falsification tests that need no ground truth, such as checking whether graphs
learned from different variable subsets agree. Those are necessary conditions
and none of them is a recovery measurement.

A generator removes the obstruction by construction. It yields the graph, the
regime path, the factor trajectory and every conditional distribution exactly,
so precision, recall, structural Hamming distance, regime accuracy and
parameter error all become computable. It also yields something an observed log
cannot give at any price, which is control over the difficulty of the problem.
The generator can be interrogated about a scenario before an estimator is
pointed at it, and \cref{sec:benchmark_identifiability} makes that
interrogation part of the protocol rather than an afterthought.

The cost is equally clear. A recovery number measured on data drawn from the
model being evaluated establishes self-consistency and nothing wider. Two
design decisions bound that cost. The first is a family of scenarios generated
under deliberate violations of the modeling assumptions, so that graceful
degradation rather than exact recovery is what is measured there. The second is
running the identical evaluation code against networks released by others,
where the ground truth was fixed by someone with no stake in this model.

\subsection{The four generating families}
\label{sec:benchmark_families}

Four families span the axes the model is built around, and each one isolates a
different component by removing the others.

\paragraph{Family I, stationary and sparse.}
Every scale carries a single regime and no factor, \(M_k = 1\) and \(r_k = 0\),
with a fixed sparse lagged graph and parameters constant in time. What remains
is structure recovery and the top-down allocation, with no latent dynamics to
confound either.

\paragraph{Family II, regime switching.}
At least two regimes act at the coarsest scale and \(r_k = 0\) everywhere, with
regime-specific count levels, duration levels and allocation scores, and a
transition matrix built so that the regimes persist and occupy their states
about equally. Which further scales carry a regime is not a free choice; it is
whatever the screen of \cref{sec:benchmark_identifiability} admits. This family
tests recovery of non-stationary behavior and is the one in which a persistent
latent state is available to confound a latent-free structure criterion.

\paragraph{Family III, shared latent factor.}
A single regime and a factor of dimension at least one, \(M_k = 1\) and
\(r_k \ge 1\), with a stable first-order factor dynamic per scale and loadings
that put correlated same-slice fluctuations on every label. The lagged graph
stays sparse, so the question is whether the factor absorbs contemporaneous
co-movement without distorting recovery of the lagged structure.

\paragraph{Family IV, deliberate misspecification.}
The same structure with one generating assumption broken on purpose, in two
variants. The heavier-tail variant fixes the coarsest-scale negative binomial
dispersion at a large constant, which drives the generated counts toward the
Poisson limit and leaves the fitted dispersion chasing a value the estimator
cannot reach. The omitted-context variant enlarges the coarsest-scale context
coefficients so that the context genuinely drives the emission level, and the
fitted model has to absorb its effect elsewhere. A third setting names the
reference scenario, the same structure with nothing misspecified, which is what
makes the degradation measurable rather than merely observed.

\paragraph{A family overrides the scenario it is given.}
The family is not a label attached to a configuration; it is a projection of
that configuration onto its own subfamily, applied before any parameter is
drawn. Family I forces \(M_k = 1\) and \(r_k = 0\) at every scale. Family II
forces \(r_k = 0\), raises the coarsest scale to \(\max(2, M_K)\), and leaves
every other scale as the configuration states, so a scale asked to carry no
regime carries none. Family III forces
\(M_k = 1\) and \(r_k \gets \max(1, r_k)\). Asking family I for four regimes
therefore returns one, and asking family II for a factor returns none. The
consequence for the grid is that the regime count is a live axis only inside
family II and the factor dimension only inside family III. A grid that swept
\(M_k\) and \(r_k\) across all families would have generated many scenarios
that are byte-identical to one another and reported them as independent
evidence.

\subsection{Scenario configuration and the two grids}
\label{sec:benchmark_grid}

A scenario fixes the tuple
\begin{align}
\bigl(
K,\;
T_1,\;
(b_k)_{k=1}^{K-1},\;
\gL,\;
(M_k),\;
(r_k),\;
(p_k),\;
(D_k^\star),\;
\varpi_0,\;
\text{family},\;
\text{seed}
\bigr),
\label{eq:bench_scenario}
\end{align}
where \(b_k\) is the branching factor from scale \(k\) to scale \(k+1\),
\(p_k\) is the maximum lag searched at scale \(k\) in
\cref{eq:model_lag_values}, \(D_k^\star\) is the maximum
in-degree of the true graph there, and \(\varpi_0\) is the structural zero
probability that fixes the zero-inflation intercept, set to \(0.1\), \(0.35\)
or \(0.65\) for the low, medium and high zero-inflation levels. The epoch
counts follow from the branching,
\begin{align}
T_{k+1} = \left\lfloor T_k / b_k \right\rfloor,
\qquad k = 1,\dots,K-1,
\label{eq:bench_epochs}
\end{align}
and the wall-clock length of an epoch follows from the length of one finest
slot together with the same branching.

Two hierarchies are used, and both are aligned to a clock rather than chosen
for arithmetic convenience. The minute grid sets \(K = 3\), \(T_1 = 20{,}160\)
and branching \((10, 6)\) at a one-minute finest slot, giving \(20{,}160\),
\(2{,}016\) and \(336\) epochs at one minute, ten minutes and one hour over
fourteen days. The day grid sets \(K = 3\), \(T_1 = 8{,}640\) and branching
\((12, 24)\) at a five-minute finest slot, giving \(8{,}640\), \(720\) and
\(30\) epochs at five minutes, one hour and one day over thirty days. The
second exists so that the study repeats when the coarsest scale is a calendar
day, which changes the number of coarsest epochs available for estimation by
an order of magnitude.

Both grids fix the same generating conditions. The finest scale is constrained
to binary event starts, \(X_t^{(1,\ell)} \in \{0,1\}\), which is a modeling
choice about the labels rather than a requirement of the event-log reading. What
makes a reconstructed trace an event log is that every episode carries a start
and an end **at the resolution of the finest scale itself**, and nothing finer:
if that scale is an hour, an episode runs from one hour to another and may begin
and end in the same one. Two episodes of a label beginning in the same finest
epoch are therefore admissible, since the log does not claim to order events
inside an epoch, and no claim below the declared finest scale is supported by
anything the generator produced. The
coarsest log-mean is set from a target rate of \(0.15\) events per finest epoch
conditional on the label not being structurally absent, since fixing the
coarsest mean directly empties the finest scale once the counts have been
spread over sixty slots. The lagged coefficients and the regime contrasts carry
multipliers of \(10\) and \(3\) on their base draws, which is what places them
in the regime \cref{sec:benchmark_identifiability} certifies. The searched lag
bound is \((p_k) = (3,2,1)\) from finest to coarsest and the true in-degree
bound is \((D_k^\star) = (0,2,1)\), so the true edges live at the ten-minute
and the hourly scales and the finest scale is left empty. The latent states are
placed by the same rule and land differently: the regime count is
\((M_k) = (1,1,M)\) on the minute grid and \((1,M,M)\) on the day grid, since
the screen of \cref{sec:benchmark_identifiability} admits a regime at the hour
scale of the day grid and at no scale below the coarsest of the minute grid,
while the factor dimension \(r_k\) is uniform across scales because the factor
clears everywhere.

The minute grid holds twenty scenarios. Family I contributes four, from
\(L \in \{5,10\}\) crossed with low and high zero inflation. Family II
contributes eight, from \(L\) crossed with \(M \in \{2,4\}\) at the coarsest
scale and one regime below it, and the two zero inflation levels. Family III contributes eight, from
\(L\) crossed with \(r_k \in \{1,2\}\) and the two levels. The day grid adds
two at medium zero inflation, one from family II with two regimes at the hour
and day scales and one from family III with \(r_k = 1\), for twenty-two
scenarios in total, each generated
at three independent replicates.

Family IV sits outside this grid and carries its own. Degradation under
misspecification is a paired comparison against an otherwise identical
well-specified scenario, so the misspecification type has to be the swept axis
rather than one setting among twenty-two. That grid holds twelve
configurations, from the three misspecification settings crossed with
\(L \in \{5,10\}\) and \(M \in \{2,4\}\) at the coarsest scale and medium zero
inflation, on the minute clock and with the same lag bound, in-degree bound,
regime placement, rate, edge-strength and regime-strength calibration as above,
again at three replicates.

One asymmetry in the search space is deliberate. The in-degree bound the
learner searches under is \(D_k = \max(D_k^\star, 1)\) rather than
\(D_k^\star\) itself, so the finest scale, which carries no true edge, is
searched with a budget of one parent per label instead of zero. A budget of
zero would make that scale trivially correct for every method and would destroy
the precision test it exists to carry.

Each replicate is split chronologically at coarsest-scale boundaries into
training, validation and test partitions holding \(70\%\), \(15\%\) and
\(15\%\) of the coarsest epochs. Splitting on coarsest-scale boundaries is what
keeps the reconciliation identity
\(\sum_{t \in \child_k(c)} X_t^{(k,\ell)} = X_c^{(k+1,\ell)}\) true inside each
partition; a split at an arbitrary finest-scale index would cut a block in half
and break it. The split is chronological rather than random because the model
is sequential and the question is generalization forward in time. Every method
receives the same training partition, the same initialization and the same
search space, namely the scenario's own \((p_k, D_k)\), since a search space
narrower than the generating one caps recall before a single score is computed.

\subsection{The identifiability screen}
\label{sec:benchmark_identifiability}

A recovery metric carries information only when the generated data contain what
is being recovered. A scenario in which a true edge leaves less evidence than
the best non-edge produces a low \(F_1\) for every consistent estimator, and
reporting that number as a property of an estimator is a category error. The
screen described here is applied to every scenario before it enters a study,
and it decided the shape of the grid in \cref{sec:benchmark_grid}.

The screen fits nothing. It generates one replicate, runs latent-state
inference at the true parameters \(\Theta^\star\) and the true graph
\(\gG^\star\), and asks what a perfect estimator would see. Write \(q\) for the
posterior that inference returns under \(\Theta^\star\) and \(\gG^\star\), and
define the unpenalized gain of a single candidate lagged parent
\(e = (\ell', d)\) for label \(\ell\) at scale \(k\) as
\begin{align}
\Delta_{k,\ell}(e)
&=
\max_{\beta,\alpha}\;
\E_q\!\left[
\log p\bigl(X^{(k,\ell)}_{1:T_k}, Y^{(k,\ell)}_{1:T_k}
\bigm| S^{(k)}, H^{(k)}, \gK^{(k)}, X^{(k+1,\ell)}, Y^{(k+1,\ell)},
\{e\}, \beta, \alpha, \Theta_k^\star\bigr)
\right]
\nonumber \\
&\qquad -
\E_q\!\left[
\log p\bigl(X^{(k,\ell)}_{1:T_k}, Y^{(k,\ell)}_{1:T_k}
\bigm| S^{(k)}, H^{(k)}, \gK^{(k)}, X^{(k+1,\ell)}, Y^{(k+1,\ell)},
\emptyset, \Theta_k^\star\bigr)
\right],
\label{eq:probe_gain}
\end{align}
with the coefficients of the candidate profiled out, exactly as the structure
step scores it. The parent totals \(X^{(k+1,\ell)}\) and \(Y^{(k+1,\ell)}\) are
in the conditioning set at a finer scale, where the likelihood is the
Dirichlet-multinomial allocation of those totals over \(\child_k(c)\), and are
absent at \(k = K\), where it is the zero-inflated negative binomial count and
negative binomial duration pair. The expectation is likewise exact at \(k = K\)
and approximate below it, since a finer-scale allocation score is evaluated at
its regime average \(\bar{\eta}_t^{(k,\ell)}\) before the likelihood is formed,
which is the substitution the structure step itself makes.

Three quantities are then read off per scale, the median of
\(\Delta_{k,\ell}\) over the true edges, the \(95\)th percentile of
\(\Delta_{k,\ell}\) over the candidates that are not true edges, and the price
\(\log T_k\) that a plain Bayesian information criterion attaches to one added
parent. That price is a full log-epoch rather than half of one because a
parent adds two coefficients, one on the count channel and one on the duration
channel, at \(\tfrac{1}{2}\log T_k\) each. The screen also reports the number
of epochs of that scale in the training split, the fraction of label-epoch
cells there carrying at least one event, and the best-permutation regime
accuracy that the true parameters themselves attain.

The scenario is worth running a structure study on only when the three gains
order correctly,
\begin{align}
\operatorname{med}_e \Delta_{k,\ell}(e)
\;\gg\;
Q_{0.95}\bigl(\Delta_{k,\ell} \mid e \notin \gP_{k,\ell}^\star\bigr)
\;>\;
\log T_k ,
\label{eq:probe_ordering}
\end{align}
and the latent metrics are interpretable only relative to the oracle regime
accuracy. \Cref{tab:probe} gives the reading for the benchmark scenario, a
family II configuration on the minute grid with five labels, two regimes and
low zero inflation.

\begin{table}[t]
\centering
\caption{Identifiability screen on the benchmark scenario, one replicate,
latent-state inference at the true parameters and the true graph. The minute
scale carries no true edge, so its true-edge gain is undefined and the entry is
left blank. Gains are in nats.}
\label{tab:probe}
\begin{tabular}{llrrrrrr}
\toprule
\(k\) & epoch & \(T_k\) (train) & density & true edges & true gain & non-edge \(Q_{0.95}\) & \(\log T_k\) \\
\midrule
1 & one minute  & 14{,}100 & 0.255 & 0 &        &  32.82 & 9.55 \\
2 & ten minutes &  1{,}410 & 0.530 & 5 &  31.50 &  11.32 & 7.25 \\
3 & one hour    &      235 & 0.856 & 4 & 182.26 & 141.91 & 5.46 \\
\bottomrule
\end{tabular}
\end{table}

The reading of \cref{tab:probe} is the design of the grid. At the ten-minute
scale a true edge is worth \(31.50\) nats against a non-edge level of \(11.32\)
and a penalty of \(7.25\), and at the hourly scale it is worth \(182.26\)
against \(141.91\) and \(5.46\), so \eqref{eq:probe_ordering} holds at both and
the edges there are recoverable in principle. The whole gain scale at the hour
sits higher, non-edges included, because an edge there is scored against the
count and duration likelihood itself rather than against an allocation of a
parent total over children. At the minute scale the ordering fails before any
estimator is involved, since a non-edge already reaches \(32.82\) nats against a
penalty of \(9.55\), so a criterion charging the plain price would accept
non-edges on evidence alone. Placing true edges there would produce recovery
numbers that measure the scenario rather than the method.

That reading is one scenario, and the placement of the edges follows from the
same screen applied to all twenty-two. Each was generated once more at a nonzero
in-degree on every scale, and the quantity recorded was the median true-edge
gain less the floor it has to clear, \(\max(Q_{0.95}, \log T_k)\). The
ten-minute scale returns \(-4.42\), \(+7.68\) and \(+15.92\) nats for families I,
II and III with \(17\) of \(22\) scenarios clearing, and the hourly scale
\(-9.89\), \(+6.29\) and \(+11.15\) with \(14\) of \(22\). The finest scale fails
in every family, at \(-6.88\), \(-16.50\) and \(-8.64\), with only \(2\) of
\(22\) clearing. Family I clears at neither of the two coarser scales on this
draw, and its four scenarios split by zero inflation rather than by anything the
screen controls: at the hourly scale the two low-inflation scenarios return
\(-37.09\) and \(-79.24\) while the two high-inflation ones return \(+17.31\) and
\(+22.21\). A denser count series gives a non-edge more to exploit there, and the signal is
not what separates those scenarios: their realized effect sizes are $0.3244$ and
$0.3407$, so the failing arm is if anything marginally the stronger one. The
coarsest coefficients are pinned to one value, $0.3$, by the bound that keeps the
count autoregression stable, and that makes the edge-strength multiplier inert
at that scale; what
it does not make constant is the difficulty, since the evidence an edge leaves is
that coefficient times the spread of the column it multiplies, which runs
$0.2655$ to $0.7907$ across the grid. The grid therefore carries true edges at the
ten-minute and the hourly scales and leaves the finest one empty.

That failure is a property of the generator rather than an accident of one draw.
The finest scale is a binary event-start process, so a lagged edge there can
only shift which slot of a block begins an episode, which is worth less than the
price of the edge and is small relative to the freedom the allocation already
has. Leaving the scale empty converts it from a failed recovery test into a
working precision test. Any edge a method reports there is structure the data do
not contain, which makes the count of such edges a direct, interpretable
false-positive measurement with no threshold to choose. The converse has to be
stated with it, so that silence there is not read as recovery. No measurement on
this grid can credit a method for finding a finest-scale edge, because none can
be placed there.

An earlier version of this grid left the coarsest scale empty as well. The
screen that justified it had been read on scenarios generated with no
coarsest-scale edge at all, where every gain it could observe there was a
non-edge gain and the whole scale of them sat below the price. Regenerating at a
nonzero in-degree and rescreening reverses that reading, and the same screen is
what keeps the finest scale out.

The latent states need the same screen, and they need it stated with a floor
rather than a ceiling alone. The ceiling is the best-permutation regime accuracy
the true parameters themselves attain. An earlier version of this section
reported that ceiling as \(0.652\) at the minute scale, \(0.661\) at ten minutes
and \(1.000\) at the hour, and concluded that a fit near either of the first two
was at the limit the data allow. That conclusion does not follow, because a
ceiling bounds nothing on its own. The generated chain spent \(0.71\) of its
epochs in one state at the minute scale, so naming that state scored \(0.71\)
and the ceiling sat below the constant predictor. A fit at that ceiling was not
a successful inference reported as a failure; it was no inference at all, and
every parameter block driven by that path was scored against noise.

The screen therefore reports both, and refuses a latent state whose ceiling does
not clear its floor. A regime is scored against the share of its modal state,
which is what a constant predictor reaches. A factor has no constant predictor,
so it is scored against the correlation a circularly shifted copy of its own
trajectory reaches, which is the null the edge confirmation already uses. Write
\(\pi^\star_{\max}\) for the modal share of the generating path at scale \(k\)
and \(a_k^\star\) for the accuracy attained under \(\Theta^\star\) and
\(\gG^\star\); the state is admissible at that scale only when
\begin{align}
a_k^\star - \pi^\star_{\max} \;\ge\; \delta ,
\label{eq:latent_ordering}
\end{align}
with \(\delta = 0.05\), so that a path beating the constant by a thousandth is
not counted as recovered.

Applied to the nine family II scenarios, \eqref{eq:latent_ordering} holds at the
coarsest scale of both grids in \(9\) of \(9\), by \(+0.468\) to \(+0.600\). It
fails at the minute scale in \(0\) of \(8\), with every margin inside \(0.021\)
of the floor, and holds at the ten-minute scale in \(1\) of \(8\), which does not
survive raising the regime contrast and is therefore a draw rather than a signal.
The hour scale of the daily grid is the exception and keeps its regime, clearing
four times out of four at \(+0.157\) and above; its blocks hold \(24\) children
against the minute grid's \(6\), and that is the channel a regime has to act
through once it no longer reaches the count emission. The contrast is not the
lever, which had to be measured before the scale could be blamed: raising it
from \(3\) to \(8\) to \(20\) moves the two finer scales from \(0\) and \(2\) of
\(9\) clearing to \(1\) and \(2\) of \(9\), while the coarsest clears \(9\) of
\(9\) throughout. The grid therefore carries \(M_k = [1, 1, M]\) on the minute
hierarchy and \([1, M, M]\) on the daily one, decided exactly as the in-degree
is.

The generating chain is built so that this comparison is worth making. Its
transition logits used to be drawn freely, so the stationary distribution was
whatever the draw implied and the modal share at the coarsest scale was
\(0.672\) at two regimes. Each row is now \(q e_m + (1 - q)\pi\), a self
transition mixed with a draw from \(\pi\), which leaves \(\pi\) stationary
exactly for any persistence \(q\), with \(\pi\) drawn within ten percent of an
equal share. The modal share fell to \(0.511\). That rebalance changed no verdict
in \eqref{eq:latent_ordering}, since the finer scales fail on the ceiling sitting
at the floor rather than on the floor being high, but it is what makes a reported
accuracy at the coarsest scale readable.

The factor clears at every scale of every scenario, by \(+0.34\) to \(+0.78\) in
correlation, including at the minute scale where no lagged edge can be placed and
no regime can be recovered. Same-slice co-movement survives the event-start
reading and persistent drift does not, which is why \(r_k\) is uniform across
scales and \(M_k\) is not. The same argument applies to every parameter block
whose identified form is coarser than its stored form.

Stated generally, the screen turns a claim about an estimator into a claim
about an estimator on a scenario that has been certified to contain the signal.
It costs one replicate and one pass of latent-state inference per scenario,
which is negligible against the cost of the studies it protects, and it is the
step that makes a low recovery number attributable.

\subsection{Methods under comparison}
\label{sec:benchmark_methods}

Seven entries share one interface, so that the protocol fits them on identical
data and the difference between two rows is the single design decision that
separates them.

The \emph{full model} runs all three learning steps at every scale. The
\emph{single-scale variant} switches off the cross-scale coupling, so no
parent-scale posterior reaches latent-state inference or parameter learning and
the regime chain at each scale sees only its own data and context. The
allocation still ties a finer scale to its parent totals, since that is what
makes the counts reconcile, so this entry isolates the value of passing the
parent posterior downward rather than the value of the hierarchy itself. The
\emph{no-factor} ablation forces \(r_k = 0\) at every scale, which sends all
same-slice co-movement into the regimes, the lagged edges or the noise. The
\emph{no-regime} ablation collapses the chain to \(M_k = 1\), which removes
non-stationarity while leaving the hierarchy and the graph intact. The
\emph{empty graph} fixes \(\gP_{k,\ell} = \emptyset\) at every scale with the
structure step switched off, leaving top-down transfer and the latent model as
the only mechanisms; it is the reference against which a structural Hamming
distance becomes meaningful, because a method whose distance exceeds it is
worse than reporting nothing. The \emph{oracle graph} installs
\(\gG^\star\) and switches the structure step off, which separates parameter
and latent-state recovery from search error.

The seventh entry, \emph{oracle parameters}, fits nothing at all. It runs
latent-state inference at \(\Theta^\star\) and \(\gG^\star\) and reports the
resulting metrics. It is a ceiling and not a competitor, and comparing a fitted
model against it answers a question no other row can, which is whether a
disappointing latent metric reflects a learner that failed or a scenario that
does not identify the quantity. The screen of
\cref{sec:benchmark_identifiability} already fixes that ceiling numerically for
the regime path, and the oracle row extends it to every metric the protocol
computes.

When a published structure learner supplies the graph, it enters through the
same interface, with the learned edges replacing \(\gG\) and nothing else
changed. One implementation detail decides whether that comparison is valid at
all. The parameter step updates the coefficient of a lagged parent only where a
coefficient for that parent already exists, so a graph installed from outside
the structure step has no effect on any likelihood until a zero coefficient is
created for each of its parents. Without that step every externally supplied
graph is scored exactly as if it had reported the empty graph, and a table of
apparently distinct methods measures nothing.

\subsection{Metrics}
\label{sec:benchmark_metrics}

\paragraph{Structure recovery, per scale.}
Write
\(E_k = \{(\ell', d, \ell) : \ell \in \gL,\; (\ell', d) \in \gP_{k,\ell}\}\)
for the lagged edge set that a collection of parent sets at scale \(k\)
defines, and
\(\hat{E}_k\) and \(E_k^\star\) for its learned and true values. The two are
compared as sets of triples, giving true positives, false positives and false
negatives and hence
\begin{align}
\mathrm{Prec}_k = \frac{|\hat{E}_k \cap E_k^\star|}{|\hat{E}_k|},
\qquad
\mathrm{Rec}_k = \frac{|\hat{E}_k \cap E_k^\star|}{|E_k^\star|},
\qquad
\mathrm{SHD}_k = \bigl|\hat{E}_k \,\triangle\, E_k^\star\bigr| ,
\label{eq:bench_structure}
\end{align}
with \(F_{1,k}\) the harmonic mean of the first two. Reporting per scale rather
than pooled is forced by the design, since one of the three scales carries no
true edges and pooling would let recall at another scale hide false positives
there. The empty-versus-empty case is scored by convention rather than by
formula. When \(\hat{E}_k\) and \(E_k^\star\) are both empty the scale is
perfectly recovered, so precision, recall and \(F_1\) are all set to \(1\);
treating them as \(0/0\) and recording zero drags the macro average down by a
whole scale in every scenario, since the finest scale is empty in all of them.
The convention has to be declared wherever the macro average over scales is
quoted, because it hands a free \(1\) at the finest scale to any method that
proposes nothing there, and a method that proposes nothing anywhere collects
that term as well.

\paragraph{Latent-state recovery.}
The regime label is identifiable only up to permutation, so the estimate
\(\hat{S}_t^{(k)} = \argmax_m \gamma_t(m)\) is scored under the relabeling that
suits it best,
\begin{align}
\mathrm{Acc}_S^{(k)}
=
\max_{\sigma \in \mathfrak{S}_{M_k}}
\frac{1}{T_k}\sum_{t=1}^{T_k}
\mathbb{I}\bigl\{\sigma(\hat{S}_t^{(k)}) = S_t^{(k)\star}\bigr\},
\label{eq:bench_accs}
\end{align}
enumerated exactly for \(M_k \le 8\) and matched greedily on the contingency
table above that. Accuracy alone rewards a degenerate estimate on an unbalanced
regime path, so the adjusted Rand index and the normalized mutual information
are reported alongside it; both are invariant to relabeling by construction and
the adjusted Rand index is corrected for chance agreement.

Two readings of \eqref{eq:bench_accs} that a table would otherwise get wrong. A
scale whose true path never switches has no regime to recover, and
\eqref{eq:bench_accs} returns exactly \(1\) there for any estimate whatsoever, so
all three figures are reported as not applicable rather than as a perfect score.
Since the grid places a regime only where the screen of
\cref{sec:benchmark_identifiability} admits one, that is the usual case at the
two finer scales of the minute hierarchy. A scale whose true path switches while
the fitted model emits a constant is scored, and the score is worth reading: the
only available relabeling maps that constant onto one true regime, so
\eqref{eq:bench_accs} returns that regime's frequency, which is what a constant
predictor earns and is a floor rather than a recovery. Regime accuracy from a
model with \(M_k = 1\) is therefore not comparable with regime accuracy from one
with \(M_k > 1\). The factor is
identified only up to a rotation, so pointwise recovery of \(H_t^{(k)}\) is not
the quantity of interest. The metric averages, over the inferred components at
a scale, the largest absolute correlation that component attains against any
true component, which is invariant to the sign and to the ordering of the
components but not to a general rotation of them.

\paragraph{Aggregation consistency.}
The model claims exact reconciliation, so the claim is measured rather than
assumed,
\begin{align}
\mathrm{RecErr}_X
=
\sum_{k=1}^{K-1}\sum_{\ell \in \gL}\sum_{c \in \gT^{(k+1)}}
\Bigl| X_c^{(k+1,\ell)} - \sum_{t \in \child_k(c)} X_t^{(k,\ell)} \Bigr|,
\label{eq:bench_recx}
\end{align}
and identically for \(\mathrm{RecErr}_Y\) with \(Y\) in place of \(X\). The
quantity is a property of whatever multiscale array it is handed. In the
protocol that array is the held-out partition of the generated data, which
satisfies the identity by construction because the finer scales were produced
by allocating the parent total, so every entry scores exactly zero and the
column certifies the generator and the split rather than separating methods.
It becomes a comparison only when it is applied to an array a method
reconstructed itself, where a scale predicted without reference to its parent
leaves a residual whose magnitude is hierarchical inconsistency in the units of
the data.

\paragraph{Held-out log evidence, and how it must be read.}
The predictive quantity is the negative log evidence on the validation
partition, \(-\sum_k \log Z_k\), where \(\log Z_k\) is what latent-state
inference returns at scale \(k\); the recovery metrics are computed on the test
partition instead. It sums a zero-inflated negative binomial count likelihood
and a negative binomial duration likelihood at the coarsest scale with a
Dirichlet-multinomial allocation likelihood at each finer scale, the allocation
terms up to an additive constant that depends only on the data, since the
block-level normalizer and the multinomial coefficient shift every candidate by
the same amount and are dropped. That constant is the reason the quantity
compares methods on one dataset and does not compare datasets. Its
magnitude also grows with \(T_k\) and with the count level, so an average of
raw values across scenarios is dominated by whichever scenario happens to be
largest. The defensible summary is the paired per-fit difference between two
methods on the same replicate, reported as the fraction of fits in which one
beats the other together with the mean difference.

\paragraph{Parameter recovery.}
Structure recovery reports which edges returned and is silent on whether the
conditional distributions attached to them are right. The parameter metric
scores the learned blocks directly against the generating ones, and it scores
only the identified form of each block. Two alignments are needed. The regime
permutation is resolved once per scale by matching identified emission levels,
and every regime-indexed block is then compared under that one relabeling. The
factor loading is identified only up to a rotation, since rotating the factors
rotates the loadings back, so the loadings are compared through their row
space, by one minus the mean best absolute cosine between the normalized true
and learned loading vectors.

Which form is identified differs by scale, and comparing raw arrays instead
charges for directions the data cannot see. At the coarsest scale the count
mean is
\(\exp\bigl(\beta_{0,\ell}^{(K)} + \beta_{S,\ell}^{(K)}[m] + \cdots\bigr)\), so
only the sum \(\beta_{0,\ell}^{(K)} + \vbeta_{S,\ell}^{(K)}\) is pinned down and
a constant may move between the two freely. At a finer scale the allocation is
a softmax over the children of one block, so every term constant inside that
block cancels exactly and only the contrast of the regime effect across regimes
survives,
\begin{align}
\text{identified level of label } \ell \text{ at scale } k < K
\;=\;
\vbeta_{S,\ell}^{(k)}
- \frac{1}{M_k}\sum_{m=1}^{M_k} \beta_{S,\ell}^{(k)}[m].
\label{eq:bench_identified}
\end{align}
The reported blocks are then the total variation distance between transition
tables, the mean absolute lag coefficient error on the true support, the
absolute learned lag mass placed off that support, the emission level error
divided by the spread of the true block, the loading subspace error and the
mean absolute difference of the child shares a regime implies for the
allocation. The scale-free entries are averaged into a single figure, and the
headline is the sample-size ladder rather than any single value, since
parameter recovery is a consistency claim and an error that does not fall with
\(T_1\) is a bias rather than an estimate.

\paragraph{Edge AUPRC.}
The candidate edge sets here are sparse, so the negative class dominates and a
single thresholded \(F_1\) confounds the quality of a ranking with the choice
of operating point. Where a method returns scored candidates rather than a
hard set, the candidates are ranked by score descending and the precision at
each recall gain is summed and divided by the number of true edges. Equally
scored candidates form one indistinguishable block and are all credited the
precision at the end of that block, so the value does not depend on the order
in which the scores happen to be enumerated, and a true edge that was never
scored contributes nothing while still counting in the denominator. The
quantity is undefined when a scale has no true edges, which is the honest
outcome rather than a value of zero.

\subsection{Protocol items that decide whether the evaluation survives review}
\label{sec:benchmark_review}

Four reporting decisions separate a structure-learning evaluation that can be
compared against published work from one that cannot. They are listed here
because each has a documented failure mode and because each is cheap to satisfy
in advance and impossible to repair afterwards.

\paragraph{Standardize, then report both sortability statistics.}
A synthetic generator that draws coefficients on a fixed positive range induces
a marginal variance ordering that agrees with the causal order, and a
continuous-optimization learner then scores well by sorting on variance rather
than by discovering anything \citep{reisach2021varsortability}. The design here
standardizes the response and the lagged design matrix before the
continuous-optimization entry is solved, which removes the raw-scale advantage.
Standardizing is necessary and not sufficient, since an \(R^2\)-based
sortability survives it \citep{reisach2023r2sortability}, so both varsortability
and \(R^2\)-sortability belong in the reported description of any generated
grid. Both are computed here. Over the $44$ scale fits of the grid that carry
true edges, two per scenario, varsortability is $0.563 \pm 0.233$ and
\(R^2\)-sortability is $0.452 \pm 0.221$, both near the chance level of $0.5$.
The generator does not induce a variance ordering that agrees with the causal
order, so no entry in the comparison tables can be scoring for that reason.

The tiled bnlearn networks used as external ground truth are the more sortable
of the two constructions. Their varsortability runs from $0.462$ on sachs to
$0.773$ on child, and their \(R^2\)-sortability from $0.621$ on child to
$1.000$ on earthquake, so a learner has ordering information there that the
generated grid does not supply. That is a property of the tiling rather than of
any method scored against it, and it is stated because the construction is
widely used and, as far as we are aware, this statistic has not been reported
for it.

\paragraph{State which structural Hamming distance is meant.}
The original definition is over partially directed acyclic graphs
\citep{tsamardinos2006mmhc}, where a reversed edge and an unoriented edge are
distinct kinds of error. Most released code computes the distance on the
directed adjacency instead, and some implementations additionally offer a
switch for whether a reversed edge costs one or two. The two definitions are
different quantities and a distance reported without saying which is not
commensurable across papers. This work computes the distance on the directed
adjacency of the lagged edge set, as in \eqref{eq:bench_structure}, counting
each missing and each spurious edge once. Within the hypothesis class searched
here every candidate edge points forward in time and its reversal is not
admissible, so the orientation ambiguity that separates the two definitions in
the static case does not arise inside the search space.

\paragraph{Report AUPRC, not AUROC alone.}
Under the class imbalance of sparse edge scoring a large number of false
positives moves the area under the receiver operating characteristic very
little while moving precision a great deal, so the two curves rank methods
differently and only the precision-recall view reflects the imbalance
\citep{davis2006prroc}. The area under the precision-recall curve is therefore
the ranking metric reported wherever a method exposes scored candidates, and
the thresholded triple of \eqref{eq:bench_structure} is reported alongside it
rather than in place of it.

\paragraph{Report regime recovery separately from edge recovery.}
A persistent latent state raises the level of every label it drives for as long
as it lasts, so each label is autocorrelated even with an empty parent set, and
a criterion that conditions on nothing attributes that autocorrelation to the
only mechanism in its hypothesis space, a lagged self edge. This is the
published mechanism by which autocorrelation inflates false positives in
time-series structure learning \citep{runge2018pitfalls}, and it is exactly the
confound the empty finest scale of \cref{sec:benchmark_grid} is built to expose.
Because a method can track the regime path well and still over-select edges, or
the reverse, the regime metrics of \eqref{eq:bench_accs} are reported per scale
and are never folded into the structural columns. Pooling the two into a single
score would hide precisely the trade-off the benchmark exists to measure.
